\documentclass{article}
\usepackage[utf8]{inputenc}
\usepackage{amsmath}
\usepackage{amsfonts}
\usepackage{amssymb}
\usepackage{graphicx}
\usepackage[margin=1in]{geometry}
\usepackage{hyperref}
\usepackage{listings}
\usepackage{xcolor}

\newcommand{\learninguniverse}[1]{}
\newcommand{\track}[1]{}
\newcommand{\module}[1]{}
\newcommand{\lesson}[1]{}
\newcommand{\overview}[1]{}
\newcommand{\overviewmarkdown}[1]{}
\newcommand{\practice}[1]{}
\newcommand{\quiz}[1]{}
\newcommand{\checkpoint}[1]{}
\newcommand{\discussion}[1]{}
\newcommand{\resource}[1]{}
\newcommand{\theory}[1]{}
\newcommand{\note}[1]{}
\newcommand{\tip}[1]{}
\newcommand{\summary}[1]{}

\title{Artificial Intelligence Notes}
\author{GATEHUB}
\date{\today}

\begin{document}

\maketitle

\learninguniverse{
title={Artificial Intelligence Notes},
description={W3Schools-style AI notes},
difficulty={Beginner to Intermediate},
estimatedHours={25},
skills={AI,Machine Learning,Python},
category={Artificial Intelligence}
}

\track{
title={Artificial Intelligence Foundations},
description={Learn AI step by step},

\module{
title={Module 1: Introduction to AI},
description={What AI is and how it is used},
estimatedHours={5},

\lesson{
title={What is Artificial Intelligence?},
\overviewmarkdown={AI is the simulation of human intelligence in machines.},
\practice{
language={python},
startercode={
num = 7
if num % 2 == 0:
    print("Even")
else:
    print("Odd")
},
expectedoutput={Odd}
},
\quiz{
question={What is AI?},
optionA={Robot},
optionB={Simulation of human intelligence},
optionC={Database},
optionD={Language},
correct={B},
explanation={AI simulates human intelligence.}
},
\checkpoint{title={Lesson 1 complete}}
}

\lesson{
title={Types of AI},
\overviewmarkdown={Narrow AI is most common today.},
\quiz{
question={Most common AI type?},
optionA={General AI},
optionB={Super AI},
optionC={Narrow AI},
optionD={Self-Aware AI},
correct={C},
explanation={Narrow AI handles specific tasks.}
},
\checkpoint{title={Module 1 complete}}
}
}

\module{
title={Module 2: Machine Learning},
description={How machines learn},
estimatedHours={8},

\lesson{
title={What is Machine Learning?},
\overviewmarkdown={ML learns from data. Accuracy = 92% in examples.},
\practice{
language={python},
startercode={
correct = 45
total = 50
print((correct / total) * 100)
},
expectedoutput={90.0}
},
\quiz{
question={Uses labeled data?},
optionA={Unsupervised},
optionB={Supervised},
optionC={None},
optionD={Reinforcement only},
correct={B},
explanation={Supervised learning uses labels.}
}
}

\lesson{
title={Train and Test},
\overviewmarkdown={Always test on unseen data.},
\discussion{prompt={Why use a test set?}}
}
}

\module{
title={Module 3: Neural Networks},
description={Deep learning basics},
estimatedHours={6},

\lesson{
title={Neurons and Layers},
\overviewmarkdown={Deep networks have many hidden layers.},
\quiz{
question={What makes a network deep?},
optionA={Many hidden layers},
optionB={GPU only},
optionC={No output layer},
optionD={One input},
correct={A},
explanation={Depth means many hidden layers.}
},
\summary{title={Course Summary},content={You learned AI and ML basics.}},
\resource{type={website},title={Python Tutorial},url={https://docs.python.org/3/tutorial/}},
\checkpoint{title={Course complete!}}
}
}

}

\end{document}
